\section{Web Interface}

The web interface is a \textbf{Flask} application (Python) that exposes the five operations to end users
via an HTML frontend. It communicates with Oracle Database through the \texttt{oracledb} Python library
by calling the stored procedures and functions defined in Section~3.

\subsection{Backend -- \texttt{app.py}}

\begin{lstlisting}
# app.py - SageFit Sports Equipment Booking System
from flask import Flask, render_template, request, redirect, url_for, flash
import oracledb, datetime

app = Flask(__name__)
app.secret_key = 'supersecretkey'

DB_USER = 'System'; DB_PASSWORD = 'password123'; DB_SID = 'localhost:1521/xe'

def get_db_connection():
    try:
        return oracledb.connect(user=DB_USER, password=DB_PASSWORD, dsn=DB_SID)
    except Exception as e:
        print("DB connection error:", e); return None

@app.route('/')
def index():
    return render_template('index.html')

# Operation 1: Register a new customer
@app.route('/register_customer', methods=['GET', 'POST'])
def register_customer():
    if request.method == 'POST':
        conn = get_db_connection()
        try:
            cur = conn.cursor()
            cur.callproc("proc_register_customer", [
                request.form['cust_id'], request.form['name'],
                request.form['cust_type'], request.form['email'],
                request.form['phone'],    request.form['address']])
            conn.commit(); flash("Customer registered successfully", "success")
        except oracledb.DatabaseError as e:
            flash(f"Error: {e}", "danger")
        finally: cur.close(); conn.close()
        return redirect(url_for('index'))
    return render_template('register_customer.html')

# Operation 2: Add a new booking
@app.route('/add_booking', methods=['GET', 'POST'])
def add_booking():
    if request.method == 'POST':
        booking_date = datetime.datetime.strptime(
            request.form['booking_date'], '%Y-%m-%d').date()
        conn = get_db_connection()
        try:
            cur = conn.cursor()
            cur.callproc("proc_add_booking", [
                request.form['booking_id'], request.form['btype'],
                booking_date, int(request.form['duration']),
                float(request.form['cost']), request.form['method'],
                request.form['contract_type'], request.form['team_code'],
                request.form['customer_id'], request.form['location_id']])
            conn.commit(); flash("Booking added successfully", "success")
        except Exception as e:
            flash(f"Error: {e}", "danger")
        finally: cur.close(); conn.close()
        return redirect(url_for('index'))
    return render_template('add_booking.html')

# Operation 3: Add a new team
@app.route('/add_team', methods=['GET', 'POST'])
def add_team():
    if request.method == 'POST':
        conn = get_db_connection()
        try:
            cur = conn.cursor()
            cur.callproc("proc_add_team", [
                request.form['team_code'], request.form['team_name'],
                int(request.form['max_members'])])
            conn.commit(); flash("Team added successfully", "success")
        except Exception as e:
            flash(f"Error: {e}", "danger")
        finally: cur.close(); conn.close()
        return redirect(url_for('index'))
    return render_template('add_team.html')

# Operation 4: Add a new member
@app.route('/add_member', methods=['GET', 'POST'])
def add_member():
    if request.method == 'POST':
        conn = get_db_connection()
        try:
            cur = conn.cursor()
            cur.callproc("proc_add_member", [
                request.form['member_id'], request.form['name'],
                request.form['surname'],   request.form['phone'],
                request.form['email'],     request.form['team_code']])
            conn.commit(); flash("Member added successfully", "success")
        except Exception as e:
            flash(f"Error: {e}", "danger")
        finally: cur.close(); conn.close()
        return redirect(url_for('index'))
    return render_template('add_member.html')

# Operation 5: View all members of a team
@app.route('/team_members', methods=['GET', 'POST'])
def team_members():
    members = []; team_code = None
    if request.method == 'POST':
        team_code = request.form.get('team_code')
        conn = get_db_connection()
        try:
            cur = conn.cursor()
            out_cursor = cur.var(oracledb.CURSOR)
            cur.callproc("proc_get_team_members", [team_code, out_cursor])
            members = out_cursor.getvalue().fetchall()
        except Exception as e:
            flash(f"Error: {e}", "danger")
        finally: cur.close(); conn.close()
    return render_template('team_members.html',
                           members=members, team_code=team_code)

if __name__ == '__main__':
    app.run(debug=True)
\end{lstlisting}

\subsection{Frontend}

\begin{figure}[H]
    \centering
    \includegraphics[width=\textwidth]{images/HomePage.png}
    \caption{Homepage -- links to all five operations}
\end{figure}

\begin{figure}[H]
    \centering
    \includegraphics[width=\textwidth]{images/Op1.png}
    \caption{Operation 1 -- Register a new customer}
\end{figure}

\begin{figure}[H]
    \centering
    \includegraphics[width=\textwidth]{images/Op2.png}
    \caption{Operation 2 -- Add a new booking}
\end{figure}

\begin{figure}[H]
    \centering
    \includegraphics[width=\textwidth]{images/Op3.png}
    \caption{Operation 3 -- Add a new team}
\end{figure}

\begin{figure}[H]
    \centering
    \includegraphics[width=\textwidth]{images/Op4.png}
    \caption{Operation 4 -- Add a new member}
\end{figure}

\begin{figure}[H]
    \centering
    \includegraphics[width=\textwidth]{images/Op5.png}
    \caption{Operation 5 -- View members of a specific team}
\end{figure}
